\makeabstract
    {
    The Role of m6A mRNA Methylation in D2 Striatal Neurons During Pavlovian and Instrumental Learning and Extinction
    }
    {
    Erica Li\textsuperscript{2}, Zhuoyue Shi\textsuperscript{1}, Nabilah Sammudin\textsuperscript{1}, Xiaoxi Zhuang\textsuperscript{1}
    }
    {
    \textsuperscript{1} Department of Neurobiology, The Neuroscience Institute, University of Chicago, Chicago, IL\\
\textsuperscript{2} Neuroscience Program, Amherst College, Amherst, MA
    }
    {
    To understand learning, it is imperative to unravel the cellular and molecular mechanisms regulating changes in protein synthesis in neurons undergoing plasticity. The post-transcriptional methylation of adenosine into N6-methyladenosine (m6A) within mRNA transcripts is particularly efficient in modifying protein synthesis in neurons. This methylation is catalyzed by the METTL14-METTL3 heterodimer complex, and previous work has found that the expression of METTL14 in D1 and D2 medium spiny neurons in the striatum---a brain region in the basal ganglia crucial for learning---plays different roles in motor learning and drug-dependent learning. The current study investigates the cell-type dependent role of METTL14 in associative learning by examining the effect of Mettl14 deletion in D2 striatal neurons on Pavlovian and operant learning and extinction. 
    }
    {
    Five A2A-Mettl14 knockout mice and 5 control mice were used in behavioral experiments carried out in Med-PC IV operant boxes. Prior to the start of each experiment, mice were food deprived for one day. For Pavlovian learning, mice first underwent 13 days of acquisition, during which they experienced a signal light followed by a sugar pellet reward. They subsequently underwent 5 days of extinction, during which the reward was removed, and 1 day of probe, during which the reward was restored. For operant conditioning, mice underwent 7 days of acquisition, during which they lever pressed for sugar pellets on a Fixed Ratio 5 schedule. This was followed by 5 days of extinction and 1 day of probe. All data analysis was performed using GraphPad Prism.
    }
    {
    No significant overall difference in head entries was observed between control and knockout mice during Pavlovian acquisition, extinction, or probe. Notably, however, knockout mice showed significantly higher baseline head entries than controls specifically on days 5 and 6 of acquisition (p \ensuremath{<} 0.05). No significant difference in lever presses was observed between control and knockout mice during operant acquisition or probe; however, knockout mice showed significantly lower lever presses on the first day of extinction compared to controls (p \ensuremath{<} 0.01), though this difference did not last for the rest of the extinction period.
    }
    {
    In general, the expression of METTL14 in D2 striatal neurons does not seem to play an important role in associative learning or extinction in mice. However, METTL14 expression in D2 striatal neurons could potentially be involved in the early refinement of behavioral patterns during associative learning, as suggested by heightened baseline head entries exhibited by knockout mice during Pavlovian acquisition.
    }

    \newpage
    
